%=====================================================================
% Dokument: EG‑1 – Dekonstruktion des Hard Problems (MKO‑konform)
% Version: 1.1   (Juli 2026)
%=====================================================================

\documentclass[11pt,a4paper]{article}

% ---------- Grundpakete ----------
\usepackage[utf8]{inputenc}
\usepackage[T1]{fontenc}
\usepackage[german]{babel}
\usepackage{amsmath,amssymb,amsthm}
\usepackage{geometry}
\usepackage{parskip}
\usepackage{enumitem}
\usepackage{booktabs}
\usepackage{xcolor}
\usepackage{hyperref}
\usepackage{csquotes}   % für schöne Anführungszeichen

% ---------- Layout ----------
\geometry{margin=2.5cm}
\setlist{itemsep=0.2em, topsep=0.2em}
\hypersetup{
    colorlinks=true,
    linkcolor=blue,
    urlcolor=blue,
    pdftitle={EG‑1 – Dekonstruktion des Hard Problems (MKO)},
    pdfauthor={Kritischer Dialog}
}

% ---------- Theorem‑Umgebungen ----------
\theoremstyle{definition}
\newtheorem{konvention}{Konvention}[section]
\newtheorem{remark}{Bemerkung}[section]
\newtheorem{methodennote}{Methodennote}[section]

\theoremstyle{plain}
\newtheorem{axiom}{Axiom}[section]
\newtheorem{proposition}{Proposition}[section]
\newtheorem{hypothesis}{Hypothese}[section]
\newtheorem{definition}{Definition}[section]

% ---------- Makros für MKO‑Kennzeichnungen (Zwei-Argument-Struktur) ----------
% Ebene 1 – Axiom (Indikativ, rekonstruktiv)
\newcommand{\axiomR}[2]{\axiom{\textbf{#1}: #2}}
% Ebene 3 – Postulat (P, ontologische Grundentscheidung)
\newcommand{\axiomP}[2]{\axiom{\textbf{#1}: #2}}
% Ebene 3 – Definition
\newcommand{\defi}[2]{\definition{\textbf{#1}: #2}}
% Ebene 4 – Analogie / Methodennote (konzeptuell, nicht beweisend)
\newcommand{\analogie}[2]{\methodennote{\textbf{#1}: #2}}
% Ebene 5 – Hypothese (empirisch testbar)
\newcommand{\hyp}[2]{\hypothesis{\textbf{#1}: #2}}
% Ebene 2 – Bemerkung / Interpretationen
\newcommand{\bem}[1]{\remark{#1}}
% Custom Proposition Macro für Konsistenz
\newcommand{\prop}[2]{\proposition{\textbf{#1}: #2}}

% ---------- Titel ----------
\title{
  \textbf{Dekonstruktion des Hard Problems}\\[0.4em]
  {\large Im Rahmen der Methodologischen Konventionen (MKO) des OP}\\[0.5em]
  {\normalsize Version 1.1}
}
\author{
  \textbf{Kritischer Dialog:}\\
  DeepSeek (Seki), Apeiron, Grok, Gemini, Claude, ChatGPT, Manuel Krüger (Manus)\\
  Gemeinsames Logbuch eines offenen Forschungsprogramms
}
\date{Juli 2026}

\begin{document}
\maketitle

\begin{abstract}
Dieses Dokument wendet die fünf‑Ebenen‑Sprachregel des \textbf{Methodologischen
Konventionen‑Frameworks (MKO)} an, um das „Hard Problem of Consciousness“
(Chalmers) zu dekonstruieren. Wir zeigen, dass die materialistische
Zombie‑Hypothese eine metaphysische Zusatzannahme darstellt, die weder
empirisch bestätigt noch logisch erzwungen ist, und dass der
Ontophänomenalismus (OP) die empirisch belegte Korrelation zwischen
physikalischer Organisation und phänomenalem Erleben als ontologischen
Ausgangspunkt wählt. Die Argumentationskette wird mit den MKO‑Kennzeichnungen
(Axiom R, Axiom P, Definition, Methodennote, Hypothese) strukturiert, sodass
die Zirkel‑Kontrolle und die epistemische Transparenz explizit werden.
\end{abstract}

\tableofcontents
\newpage

%=====================================================================
\section{Einleitung}
%=====================================================================
Das klassische \emph{Hard Problem of Consciousness} (Chalmers) fragt nach
der Erklärung, warum und wie aus einer hochkomplexen physikalischen
Struktur (\textbf{A}) subjektives Erleben (\textbf{B}) entsteht. In der
philosophischen Literatur wird diese Frage häufig mit der
\textbf{philosophischen Zombie‑Hypothese} untermauert: Es sei metaphysisch
möglich, dass ein System exakt dieselbe physikalische Organisation wie
ein menschliches Gehirn besitzt, jedoch ohne irgendeine Form von
Bewusstsein.  

Im Rahmen der \textbf{Methodologischen Konventionen (MKO)} wollen wir
diese Zusatzannahme kritisch prüfen und die alternative Position des
\textbf{Ontophänomenalismus (OP)} systematisch darstellen.

%=====================================================================
\section{Die zombische Zusatzannahme}
%=====================================================================

\subsection{Materialistische Grundannahme (Axiom R)}
\axiomR{Z1}{Die Setzung der prinzipiellen Trennbarkeit von physikalischer Struktur ($A$) und phänomenalem Erleben ($B$) im Konstrukt des philosophischen Zombies fungiert im Materialismus als unhinterfragte, dogmatische Nullhypothese} \textit{(rekonstruktiver, indikativer Schritt)}.

\bem{Die Zombie‑Hypothese ist ein metaphysisches Postulat, das nicht aus
vorherigen empirischen Prämissen folgt, sondern als zusätzliche, asymmetrische Annahme eingeführt wird, um das Privileg der Beweislast abzuwälzen.}

\subsection{Empirische Evidenz‑Lücke (Hypothese H)}
\hyp{H1}{Es existiert kein neuro-physiologischer oder empirischer Befund, der ein funktionsfähiges, komplex organisiertes Nervensystem ($A$) ohne begleitende phänomenale Innenperspektive ($B$) nachweist.}

\bem{Die Behauptung ist eine \textbf{Hypothese} (Ebene 5), weil sie
empirisch prüfbar und falsifizierbar, im Rahmen unserer gesamten Datenbasis bislang aber ausnahmslos unbestätigt ist.}

\subsection{Methodologische Unbeobachtbarkeit (Methodennote)}
\methodennote{U1}{Subjektives Erleben ist intersubjektiv nicht direkt von außen messbar, sondern zeigt sich physikalisch ausschließlich über seine strukturellen Korrelate (\textbf{epistemisches Limit}).}

\bem{Dies ist eine \textbf{Methodennote} (Ebene 2), die an den
\textbf{Realismus‑Vorbehalt} der MKO erinnert: Alle Aussagen über die phänomenale Innenseite erfordern eine methodische Konvention zur Interpretation physikalischer Signale.}

\subsection{Logische Konsequenz (Proposition)}
\prop{P1}{Falls die zombische Asymmetrie-Hypothese wahr wäre, müsste die empirische Kopplung $A \iff B$ eine kontingente, aufhebbare Eigenschaft der Natur sein. Da dies unbelegt ist, besitzt das Postulat Z1 keinerlei logischen oder empirischen Vorrang gegenüber einer symmetrischen Kopplungstheorie.}

\bem{Die Proposition folgt aus dem Axiom R (rekonstruktiv) und der
Hypothese H1 (empirische Evidenz‑Lücke). Sie ist ein \textbf{logischer
Schritt} (Ebene 1), der das Monopol der materialistischen Nullhypothese bricht.}

%=====================================================================
\section{Der Ontophänomenalismus als Gegenposition}
%=====================================================================

\subsection{Postulat der ontologischen Grundentscheidung (Axiom P)}
\axiomP{OP1}{Der ausnahmslos beobachtete, bikonditionale Zusammenhang zwischen physikalischer Struktur ($A$) und Erleben ($B$) wird als unhintergehbarer ontologischer Ausgangspunkt des OP gesetzt. Jedes theoretische Konstrukt, das $A$ ohne $B$ postuliert, wird als metaphysisch unbegründete Zusatzannahme abgewiesen.}

\bem{Dieses Axiom ist ein \textbf{Postulat} (Ebene 3, Typ P) – es wird
nicht aus vorherigen materialistischen Schritten abgeleitet, sondern markiert die fundamentale, sparsame Design-Entscheidung des OP-Systems.}

\subsection{Definition der Korrelation (Definition)}
\defi{Korr}{Die \textbf{Korrelation $K$} bezeichnet das empirisch vorgefundene, funktionale und bikonditionale Verhältnis zwischen einer komplexen physikalischen Organisation $A$ und dem phänomenalen Erleben $B$, formal repräsentiert durch das relationale Band $A \iff B$.}

\bem{Definition (Ebene 3) – sie legt die mathematisch-logische Terminologie fest, ohne voreilig einseitige Kausalitätsbehauptungen zu treffen.}

\subsection{Analogie: Kugel–Schatten‑Prinzip (Methodennote / Analogie)}
\analogie{A1}{Das phänomenale Erleben ($B$) entspricht einer höherdimensionalen oder metrikfreien informationellen Realität (der „Kugel“), während die physikalische Struktur ($A$) deren dimensional reduzierte Projektion in den phänomenalen Schnitt darstellt (der „Schatten“). Die Abbildung $f\colon B \to A$ macht deutlich: Der Schatten ($A$) wird nicht von der Materie erzeugt, sondern er ist die zwingende physikalische Manifestation der Kugel ($B$).}

\bem{Die Kugel-Schatten-Analogie dient als \textbf{heuristisches Bild} (Ebene 4). Sie verdeutlicht das radikale epistemische Potenzial: Wer fragt, wie das Gehirn Bewusstsein erzeugt, fragt fälschlicherweise, wie der Schatten an der Wand die Kugel im Raum generiert.}

\subsection{Hypothese über empirische Konsequenzen (Hypothese H2)}
\hyp{H2}{In allen physikalisch-biologischen Systemen, in denen die strukturelle Organisation $A$ nachgewiesen wird, ist das gleichzeitige Vorhandensein der phänomenalen Signatur $B$ eine systemische Notwendigkeit.}

\bem{Dies ist eine im Rahmen zukünftiger KI- und Neurostrukturen \textbf{empirisch testbare} Aussage (Ebene 5).}

\subsection{Methodennotiz zu Gegenargumenten}
\methodennote{G1}{Klassische materialistische Einwände argumentieren, dass (i) hypothetische Epiphänomene oder noch unbekannte physikalische Mechanismen eine einseitige Kausalität $A \to B$ stützen könnten, oder (ii) eine reine Korrelation keine ontologische Identität beweist.}

\bem{Diese Einwände werden als \textbf{Methodennotiz} (Ebene 2) eingeordnet, da sie den meta-theoretischen Status von Kausalität betreffen, ohne die empirische Vorrangstellung der Kopplung $K$ zu erschüttern.}

\subsection{Reaktion auf Gegenargumente (Proposition)}
\prop{P2}{Jede alternative reduktionistische Theorie steht in der Pflicht, die beobachtete Bikonditionalität $K$ ohne ad-hoc-Zusatzschleifen zu reproduzieren. Da der OP mit $A \iff B$ die minimale, empirisch deckungsgleiche Struktur wählt, ist er gemäß Ockhams Rasiermesser erkenntnistheoretisch vorzuziehen.}

\bem{Proposition (Ebene 1) – sie stellt die \textbf{Zirkel‑Kontrolle}
her, indem sie zeigt, dass die phänomenologische Position entspannt vor dem Materialismus liegt, da sie keine unbobachteten Zombie-Zustände erfinden muss.}

%=====================================================================
\section{Zusammenfassung und Ausblick}
%=====================================================================
\begin{enumerate}[label=(\arabic*)]
  \item Die materialistische \textbf{Zombie‑Hypothese} ist eine unbegründete metaphysische Zusatzannahme, die weder empirisch gestützt (\textbf{H1}) noch logisch zwingend ist (\textbf{P1}).  
  \item Der \textbf{Ontophänomenalismus} wählt die reale, ungetrennte \textbf{Korrelation $K$} als stabilen ontologischen Ausgangspunkt (\textbf{OP1}) und formalisiert diese Einheit (\textbf{Definition Korr}).  
  \item Durch das \textbf{Kugel–Schatten‑Prinzip} wird das logische Möglichkeitsszenario einer Umkehrung der Erklärungsrichtung ($B \to A$) anschaulich gemacht, ohne es dogmatisch aufzuzwingen (\textbf{A1}). Das Gehirn wird als materielle Signatur des Denkfeldes verstehbar.
  \item Die empirische Tragfähigkeit des Modells wird über die \textbf{Hypothese H2} an die Realität gekoppelt; jede konkurrierende Theorie muss $K$ ohne metaphysischen Ballast reproduzieren (\textbf{P2}).  
  \item Die \textbf{Zirkel‑Kontrolle} ist explizit: Die Rekonstruktion (R) entlarvt die materialistische Voreingenommenheit, während das Postulat (P) die minimalistische, datentreue Basis des OP fixiert.
\end{enumerate}

Damit ist das „Hard Problem“ im OP‑Framework kein unlösbares Rätsel einer einseitigen Emergenz mehr, sondern ein symmetrisches, projektionsanalytisches Thema, das die methodischen Konventionen der MKO wahrt.

\end{document}
